% ============================================================
%  Section 1 — How to Practice? / 如何练习
%  原书 "How to Practice?" 段落
% ============================================================
\chapter*{How to Practice? / 如何练习}
\addcontentsline{toc}{chapter}{How to Practice? / 如何练习}

\eng{Before we dive into the core math rock guitar techniques, I would like
to share an effective practice strategy. This approach is aimed at learning
new techniques or refining existing ones. It may sound counterintuitive,
but it involves intentionally making mistakes. This process of generating
errors is essential for skill improvement. Here is how it works:}

\chn{在深入学习核心数学摇滚吉他技巧之前，我想分享一个高效的练习策略。
这个方法旨在学习新技术或打磨已有技巧。听起来可能反直觉，
但它需要你有意地制造错误——产生错误的过程是提升技能的关键。具体方法如下：}

\vspace{0.5cm}

\eng{\textbf{1. Generate Errors:} The brain processes errors after each practice
session. When you return to practice, you will begin to see improvements in
that skill. Therefore, aim to perform as many repetitions as possible of the
technique you are practicing within your dedicated practice time.}

\chn{\textbf{1. 制造错误：}大脑会在每次练习后处理错误信息。
当你再次练习时，你会发现在该项技能上有所进步。
因此，在练习时间内，尽量多做你正在练习的技巧的重复。}

\eng{\textbf{2. Use a Metronome:} A metronome should be your ally in this
process. Start with a tempo where you can comfortably play the exercise,
then gradually increase it (typically by 3 bpm increments) until you reach a
point where you are no longer comfortable. This is the sweet spot for
improvement because you will be generating more errors in your repetitions,
pushing your brain's neural plasticity to adapt to this new challenge.}

\chn{\textbf{2. 使用节拍器：}节拍器应该成为你的好帮手。
从一个你能舒适演奏练习的速度开始，然后逐渐加快（通常每次增加 3 BPM），
直到你感到不再舒适为止。这就是进步的最佳点——因为你在重复中产生了更多错误，
推动大脑的神经可塑性去适应新的挑战。}

\eng{\textbf{3. Focused Practice:} When practicing a technique (skill), it is
crucial to have a specific focus, even if it is just one aspect of the
technique. For instance, you might dedicate a practice session to observing
the movement of one of your hands or focusing on the positioning of one of
your hands. By rotating your focus in different practice sessions, you can
systematically improve various aspects of your technique.}

\chn{\textbf{3. 专注练习：}练习一项技巧时，有一个明确的专注点至关重要，
哪怕只是该技巧的一个方面。例如，你可以用一次练习专门观察某只手
的运动，或专注于某只手的定位。通过在每次练习中轮换关注点，
你可以系统性地提升技巧的各个方面。}

\vspace{0.3cm}

\eng{In summary, I encourage you to try this practice strategy for the
upcoming techniques and exercises. After each practice session, take a few
minutes to allow your brain to simply reflect on what you have practiced.
Good luck and may your practice sessions be productive and rewarding!}

\chn{总之，我鼓励你在接下来的技巧和练习中尝试这个策略。
每次练习后，花几分钟让大脑沉淀一下刚刚练习的内容。
祝你好运，愿你的每次练习都高效而充实！}
